\subsubsection{Hahn-Banach 定理}\label{subsec:hahn-banach-theorem}

在前一节的最后一个例子中, 用到了以下结果.

\begin{Theorem}
\label{thm:hahn-banach-extension}
设 $X$ 为 Banach 空间, $M$ 为 $X$ 的线性子空间. 设 $l_M$ 为从 $M$ 到 $\mathbb{R}$ ($\mathbb{C}$) 的有界线性映射. 那么存在线性泛函 $l\in X^*$, 使 $l|_M=l_M$ 且 $\|l\|=\|l_M\|$.
\end{Theorem}

换言之, 定义在任意线性子空间上的线性泛函, 都可以在不增加范数的情况下延拓. 这一结果以及可推出它的各种更一般的结果, 称为 Hahn-Banach 定理. 一些更一般的版本与凸分析和优化有关; 由于本书不深入讨论这些主题, 就不陈述那些结果. Hahn-Banach 定理保证了许多线性泛函的存在性, 因为总可以在任意有限维子空间上定义线性泛函, 再将它们延拓. 特别地, 给定任意非零 $x\in X$, 存在 $f\in X^*$, 使 $f(x)\ne0$.

Theorem~\ref{thm:hahn-banach-extension} 的证明请参阅泛函分析教材. 但对 Hilbert 空间的情形, 这一结果几乎显然. 先由连续性将 $l$ 延拓到 $M$ 的闭包, 在 $M^\perp$ 上令 $l=0$, 再由线性性在整个空间上定义 $l$. 留给读者验证, 这样定义的 $l$ 具有定理所声称的性质; 见 Problem~\ref{ex:hahn-banach-hilbert-extension-details}.
